% !PRINTED 90
\weberpage{90}
%<<CONT

Herewith, however, everything is exhausted that we must require of $m$. For
when this condition f) is satisfied, we need only assume, e.\,g., the
$e_{1,h_1}, e_{2,h_2}, \dots e_{\varrho,h_\varrho}$ not divisible by $p$,
and the remaining $e$ divisible by $p$; then the determinant

\[
\begin{vmatrix}
e_{1,h_1}\,i_{1,h_1}, & \dots & e_{1,h_\varrho}\,i_{1,h_\varrho} \\
\cdot\ \cdot\ \cdot & \cdots & \cdot\ \cdot\ \cdot \\
e_{\varrho,h_1}\,i_{\varrho,h_1}, & \dots & e_{\varrho,h_\varrho}\,i_{\varrho,h_\varrho}
\end{vmatrix}
\equiv e_{1,h_1}\,e_{2,h_2} \dots e_{\varrho,h_\varrho} \pmod{p}
\]

is therefore certainly not divisible by $p$.

And under this assumption the condition 6. can also still be satisfied.
To see this, let $p^\pi$ denote any one of the invariants, say $i_1$, and
let $c_h$ denote the index modulus which by f) is divisible by $p^\pi$,
so that

\weq{25}{c_h = p^\pi c_{1,h}}

holds. If then $q_h$ is the prime divisor of $m$ belonging to the index
modulus $c_h$, then $p$ is either equal to $q_h$ or a divisor of $q_h - 1$.
By the assumption last made, $e_{1,h}$ is not divisible by $p$, while
$e_{2,h}, e_{3,h}, \dots e_{\nu,h}$ are divisible by $p$. The products

\[
c_{1,h}\,e_{1,h}, \quad c_{2,h}\,e_{2,h}, \; \dots \; c_{\nu,h}\,e_{\nu,h}
\]

are then not all to be divisible by $q_h$ or by $q_h - 1$.

If, first, $q_h$ is a simple divisor of $m$, then $c_h = q_h -
1$, and by (25) $c_{1,h}$ is not divisible by $q_h - 1$. One can
therefore still assume $e_{1,h}$ in such a way that $c_{1,h}\,e_{1,h}$ is not divisible by
$q_h - 1$. But if $q_h$ is a $\varkappa$-fold divisor of
$m$ and $\varkappa > 1$, and if $q_h = p$, then, when $p$ is odd,
$c_h = p^{\varkappa - 1}(p-1)$, and by the hypothesis b)
$\pi \leqq \varkappa - 1$. If then $c_h$ is to be divisible by $p^\pi$,
we must have $\pi = \varkappa - 1$, and (25) gives $c_{1,h} =
p - 1$. But if $p = 2$, then $c_h = 2$ or $= 2^{\varkappa - 1}$,
and hence by the hypothesis d) we must have $\pi = 1$ or $= \varkappa -
2$, and $c_{1,h} = 1$. One can therefore
% [sic: Weber prints pi <= kappa-1; hypothesis b) gives pi >= kappa-1, which is what makes divisibility force equality.]
% [sic: Weber prints c_h=2^{kappa-1} in the second 2-adic case; the index modulus is 2^{kappa-2}, as stated in Theorem 4, and that value is required for c_{1,h}=1.]

\noindent\rule{4cm}{0.4pt}

\footnotesize
\gesperrt{Dirichlet} with application of the integral calculus. (Abhandlungen
der Berliner Akademie 1837; \gesperrt{Dirichlet's Werke}, Bd.\,I, Nr.\,21;
\gesperrt{Dirichlet-Dedekind}, Vorlesungen \"uber Zahlentheorie, Supplement VI;
\gesperrt{Bachmann}, Analytische Zahlentheorie, Abschn.\,IV). If we retain the
notation used above, then, in order to satisfy the requirement f) for a
given $p$, one may take up $p^{\pi+1}$ as a factor in $m$. But then,
if $\varrho > 1$, a further $\varrho - 1$ prime factors
$q$ must occur in $m$, which are bound to congruences of the form $q \equiv 1
\pmod{i_1} \dots$.
\normalsize
%CONT>>
